\renewcommand{\tema}{Lección: Dinámica}

\section{Repaso de vectores}

\begin{seccion}
Antes de estudiar las fuerzas, es fundamental comprender las herramientas matemáticas que utilizaremos. Las fuerzas son magnitudes vectoriales, por lo que necesitamos dominar operaciones con vectores y trigonometría básica.
\end{seccion}

\subsection{Vectores: Concepto y Representación}

\begin{seccion}
\begin{definicion}
    \textbf{Vector:} Magnitud que posee tanto cantidad (módulo o magnitud) como dirección y sentido.
\end{definicion}

\textbf{Ejemplos de magnitudes vectoriales:}
\begin{itemize}
    \item Fuerza: 20 N hacia el norte
    \item Velocidad: 30 m/s a 45$^\circ$ sobre la horizontal
    \item Desplazamiento: 10 m hacia arriba
\end{itemize}

\textbf{Magnitudes escalares} (solo tienen magnitud):
\begin{itemize}
    \item Masa: 5 kg
    \item Tiempo: 10 s
    \item Temperatura: 25 $^\circ$C
\end{itemize}

\vspace{0.3cm}

\textbf{Representación gráfica:} Un vector se representa como una flecha donde:
\begin{itemize}
    \item La longitud representa la magnitud
    \item La dirección de la flecha indica la dirección y sentido
    \item El punto de aplicación indica dónde actúa
\end{itemize}
\end{seccion}

\subsubsection{Componentes Rectangulares de un Vector}

\begin{seccion}
Todo vector en el plano puede descomponerse en dos componentes perpendiculares (horizontal y vertical):

\vspace{0.3cm}

\begin{ecuacion}
\textbf{Componentes de un vector:}

Para un vector $\vec{F}$ con magnitud $F$ y ángulo $\theta$ respecto al eje horizontal:

$$F_x = F \cos\theta \quad \text{(componente horizontal)}$$
$$F_y = F \sin\theta \quad \text{(componente vertical)}$$

La magnitud se calcula con:
$$F = \sqrt{F_x^2 + F_y^2}$$

El ángulo se obtiene con:
$$\theta = \arctan\left(\frac{F_y}{F_x}\right)$$
\end{ecuacion}

\begin{nota}
    \textbf{Nota importante:} El ángulo $\theta$ se mide desde el eje $x$ positivo (horizontal hacia la derecha) en sentido antihorario.
\end{nota}

\end{seccion}

\subsubsection{Diagrama: Descomposición Vectorial}

\begin{center}

\begin{tikzpicture}[scale=1.2]
    % Ejes
    \draw[thick, ->] (0,0) -- (5,0) node[right] {$x$};
    \draw[thick, ->] (0,0) -- (0,4) node[above] {$y$};
    
    % Vector principal
    \draw[very thick, ->, color=azuloscuro] (0,0) -- (4,3) node[midway, above left] {$\vec{F}$};
    
    % Componentes
    \draw[thick, ->, color=verdecorrecto, dashed] (0,0) -- (4,0) node[midway, below] {$F_x = F\cos\theta$};
    \draw[thick, ->, color=verdecorrecto, dashed] (4,0) -- (4,3) node[midway, right] {$F_y = F\sin\theta$};
    
    % Ángulo
    \draw[color=azulmedio] (1,0) arc (0:36.87:1);
    \node[color=azulmedio] at (1.3,0.3) {$\theta$};
    
    % Líneas punteadas auxiliares
    \draw[dotted] (0,3) -- (4,3);
\end{tikzpicture}
\shorthandon{>}\shorthandon{<}
\end{center}

\subsection{Tabla Trigonométrica de Valores Frecuentes}

\begin{seccion}
En problemas de mecánica, ciertos ángulos aparecen con frecuencia. Es útil memorizar sus valores:

\vspace{0.5cm}

\begin{center}
\begin{tabular}{|c|c|c|c|}
\hline
\rowcolor{azulmedio!30}
\textbf{Ángulo} & \textbf{sen $\theta$} & \textbf{cos $\theta$} & \textbf{tan $\theta$} \\
\hline
0$^\circ$ & 0 & 1 & 0 \\
\hline
30$^\circ$ & $\frac{1}{2}$ = 0.5 & $\frac{\sqrt{3}}{2}$ $\approx$ 0.866 & $\frac{\sqrt{3}}{3}$ $\approx$ 0.577 \\
\hline
45$^\circ$ & $\frac{\sqrt{2}}{2}$ $\approx$ 0.707 & $\frac{\sqrt{2}}{2}$ $\approx$ 0.707 & 1 \\
\hline
60$^\circ$ & $\frac{\sqrt{3}}{2}$ $\approx$ 0.866 & $\frac{1}{2}$ = 0.5 & $\sqrt{3}$ $\approx$ 1.732 \\
\hline
90$^\circ$ & 1 & 0 & $\infty$ (indefinido) \\
\hline
\end{tabular}
\end{center}

\vspace{0.3cm}

\textbf{Observaciones importantes:}
\begin{itemize}
    \item sen 0$^\circ$ = 0, cos 0$^\circ$ = 1 (vector completamente horizontal)
    \item sen 90$^\circ$ = 1, cos 90$^\circ$ = 0 (vector completamente vertical)
    \item sen 30$^\circ$ = cos 60$^\circ$ y sen 60$^\circ$ = cos 30$^\circ$ (ángulos complementarios)
    \item A 45$^\circ$, las componentes horizontal y vertical son iguales
\end{itemize}
\end{seccion}
\newpage

\subsection{Suma Vectorial}

\begin{seccion}
\textbf{Método de componentes:} Para sumar varios vectores:

\textbf{Paso 1:} Descomponer cada vector en sus componentes $x$ y $y$

\textbf{Paso 2:} Sumar todas las componentes $x$ entre sí:
$$R_x = F_{1x} + F_{2x} + F_{3x} + ...$$

\textbf{Paso 3:} Sumar todas las componentes $y$ entre sí:
$$R_y = F_{1y} + F_{2y} + F_{3y} + ...$$

\textbf{Paso 4:} El vector resultante tiene componentes $(R_x, R_y)$ y magnitud:
$$R = \sqrt{R_x^2 + R_y^2}$$

\vspace{0.3cm}

\textbf{Ejemplo:} Si $\vec{F}_1 = 30$ N a 0$^\circ$ y $\vec{F}_2 = 40$ N a 90$^\circ$:
\begin{itemize}
    \item $F_{1x} = 30$ N, $F_{1y} = 0$ N
    \item $F_{2x} = 0$ N, $F_{2y} = 40$ N
    \item $R_x = 30$ N, $R_y = 40$ N
    \item $R = \sqrt{30^2 + 40^2} = \sqrt{900 + 1600} = \sqrt{2500} = 50$ N
\end{itemize}
\end{seccion}

\section{Concepto de Fuerza}

\begin{seccion}

\begin{definicion}
    \textbf{Fuerza:} Interacción capaz de cambiar el estado de movimiento o la forma de un cuerpo.
\end{definicion}

\textbf{Características fundamentales de la fuerza:}
\begin{enumerate}
    \item Es una magnitud vectorial (tiene magnitud, dirección y sentido)
    \item Puede producir aceleración (cambio en la velocidad)
    \item Puede deformar objetos
    \item Siempre involucra la interacción entre dos cuerpos
\end{enumerate}

\textbf{Unidad de fuerza en el SI:} Newton (N)

$$1 \text{ N} = 1 \text{ kg} \cdot \text{m/s}^2$$

\begin{definicion}
    \textit{Un Newton es la fuerza necesaria para acelerar un kilogramo a razón de un metro por segundo cada segundo.}
\end{definicion}

\end{seccion}
\newpage
\subsection{Tipos de Fuerzas}

\begin{seccion}

\begin{definicion}
    \textbf{Fuerzas de contacto:} Requieren contacto físico directo
\end{definicion}

\begin{itemize}
    \item \textbf{Fuerza normal ($N$):} Perpendicular a la superficie de contacto
    \item \textbf{Fuerza de fricción ($f$):} Paralela a la superficie, se opone al movimiento
    \item \textbf{Tensión ($T$):} Fuerza ejercida por cuerdas o cables
    \item \textbf{Fuerza elástica ($F_e$):} Ejercida por resortes deformados
\end{itemize}

\begin{definicion}
    \textbf{Fuerzas de campo:} Actúan a distancia sin contacto físico
\end{definicion}

\begin{itemize}
    \item \textbf{Peso ($W$):} Fuerza gravitacional terrestre sobre objetos
    \item \textbf{Fuerza gravitacional ($F_g$):} Atracción entre masas
    \item \textbf{Fuerza eléctrica:} Entre cargas eléctricas
    \item \textbf{Fuerza magnética:} Entre imanes o corrientes
\end{itemize}
\end{seccion}

\subsection{Carácter Vectorial y Superposición de Fuerzas}

\begin{seccion}
\begin{definicion}
    \textbf{Principio de superposición:} Cuando varias fuerzas actúan simultáneamente sobre un objeto, el efecto total es equivalente a una sola fuerza igual a la suma vectorial de todas ellas.
\end{definicion}

\begin{ecuacion}
\textbf{Fuerza neta o resultante:}
$$\vec{F}_{neta} = \vec{F}_1 + \vec{F}_2 + \vec{F}_3 + ... = \sum \vec{F}$$
\end{ecuacion}

\textbf{Método para calcular la fuerza neta:}
\begin{enumerate}
    \item Identificar todas las fuerzas que actúan sobre el objeto
    \item Descomponer cada fuerza en componentes $x$ y $y$
    \item Sumar las componentes: $F_{neta,x} = \sum F_x$ y $F_{neta,y} = \sum F_y$
    \item Calcular magnitud y dirección de la resultante
\end{enumerate}
\end{seccion}

\subsection{Diagrama de Cuerpo Libre}

\begin{seccion}
\begin{definicion}
    \textbf{Diagrama de cuerpo libre (DCL):} Representación esquemática que muestra todas las fuerzas que actúan sobre un objeto, representado como un punto.
\end{definicion}

\textbf{Pasos para construir un DCL:}
\begin{enumerate}
    \item Aislar el objeto de interés
    \item Representarlo como un punto
    \item Dibujar todas las fuerzas que actúan sobre él como vectores
    \item Etiquetar cada fuerza claramente
    \item No incluir fuerzas que el objeto ejerce sobre otros cuerpos
\end{enumerate}
\end{seccion}

\subsubsection{Ejemplo de Diagrama de Cuerpo Libre}

\begin{center}
\shorthandoff{>}\shorthandoff{<}
\begin{tikzpicture}[scale=1]
    % Objeto (bloque sobre superficie)
    \draw[fill=fondogris] (0,0) rectangle (2,1.5);
    \node at (1,0.75) {$m$};
    
    % Superficie
    \draw[thick] (-1,-0.1) -- (4,-0.1);
    \draw[pattern=north east lines] (-1,-0.1) rectangle (4,-0.3);
    
    % Fuerza aplicada
    \draw[very thick, ->, color=azuloscuro] (2,0.75) -- (4,0.75) node[above, midway] {$\vec{F}$};
    
    \node at (1, -1) {Situación física};
\end{tikzpicture}
\shorthandon{>}\shorthandon{<}

\vspace{1cm}

\shorthandoff{>}\shorthandoff{<}
\begin{tikzpicture}[scale=1.5]
    % Punto (objeto aislado)
    \filldraw[color=azulmedio] (0,0) circle (0.1);
    
    % Peso
    \draw[very thick, ->, color=azuloscuro] (0,0) -- (0,-2) node[right, midway] {$\vec{W} = m\vec{g}$};
    
    % Normal
    \draw[very thick, ->, color=verdecorrecto] (0,0) -- (0,1.5) node[right, midway] {$\vec{N}$};
    
    % Fuerza aplicada
    \draw[very thick, ->, color=azuloscuro] (0,0) -- (2,0) node[above, midway] {$\vec{F}$};
    
    % Fricción
    \draw[very thick, ->, color=red] (0,0) -- (-1.2,0) node[above, midway] {$\vec{f}$};
    
    \node at (0, -2.5) {Diagrama de Cuerpo Libre};
\end{tikzpicture}
\shorthandon{>}\shorthandon{<}
\end{center}


\section{Leyes de Newton del Movimiento}

\subsection{Primera Ley de Newton: Ley de Inercia}

\begin{seccion}
\begin{definicion}
    \textbf{Enunciado:} Un objeto permanece en reposo o en movimiento rectilíneo uniforme (velocidad constante) a menos que una fuerza neta actúe sobre él.
\end{definicion}

\begin{ecuacion}
\textbf{Condición matemática:}
$$\sum \vec{F} = 0 \quad \Rightarrow \quad \vec{v} = \text{constante}$$
\end{ecuacion}

\begin{definicion}
    \textbf{Concepto de inercia:} Tendencia de un objeto a mantener su estado de movimiento. La inercia es directamente proporcional a la masa.
\end{definicion}

\textbf{Implicaciones:}
\begin{itemize}
    \item Si $\sum \vec{F} = 0$ y el objeto está en reposo ($v = 0$), permanecerá en reposo
    \item Si $\sum \vec{F} = 0$ y el objeto se mueve, continuará con velocidad constante
    \item Se requiere una fuerza neta para cambiar la velocidad (acelerar o frenar)
\end{itemize}

\textbf{Ejemplos cotidianos:}
\begin{itemize}
    \item Un libro sobre una mesa permanece en reposo
    \item Un disco de hockey sobre hielo continúa deslizándose con poca fricción
\end{itemize}
\end{seccion}

\subsection{Segunda Ley de Newton: Ley de Fuerza y Aceleración}

\begin{seccion}
\begin{definicion}
    \textbf{Aceleración:} La aceleración de un objeto es directamente proporcional a la fuerza neta que actúa sobre él e inversamente proporcional a su masa.
\end{definicion}

\begin{ecuacion}
\textbf{Ecuación fundamental:}
$$\vec{F}_{neta} = m\vec{a}$$

o en forma de componentes:
$$F_x = ma_x \quad , \quad F_y = ma_y$$

donde:
\begin{itemize}
    \item $\vec{F}_{neta}$ = suma vectorial de todas las fuerzas (N)
    \item $m$ = masa del objeto (kg)
    \item $\vec{a}$ = aceleración resultante (m/s$^2$)
\end{itemize}
\end{ecuacion}

\textbf{Consecuencias importantes:}
\begin{enumerate}
    \item La aceleración tiene la misma dirección que la fuerza neta
    \item Si $F$ aumenta, $a$ aumenta proporcionalmente (con $m$ constante)
    \item Si $m$ aumenta, $a$ disminuye (con $F$ constante)
    \item Sin fuerza neta ($F = 0$), no hay aceleración ($a = 0$)
\end{enumerate}
\end{seccion}

\subsubsection{Masa: Medida de la Inercia}

\begin{seccion}
\begin{definicion}
    \textbf{Masa ($m$):} Cantidad de materia en un objeto y medida de su inercia.
\end{definicion}

\textbf{Características:}
\begin{itemize}
    \item Unidad SI: kilogramo (kg)
    \item Mayor masa implica mayor resistencia al cambio de movimiento
\end{itemize}

\textbf{Diferencia masa vs peso:}
\begin{center}
\begin{tabular}{|l|l|l|}
\hline
\rowcolor{azulmedio!30}
\textbf{Propiedad} & \textbf{Masa} & \textbf{Peso} \\
\hline
Naturaleza & Cantidad de materia & Fuerza gravitacional \\
\hline
Tipo & Escalar & Vector \\
\hline
Unidad & kg & N \\
\hline
Depende de ubicación & No & Sí \\
\hline
\end{tabular}
\end{center}
\end{seccion}

\newpage

\subsubsection{Peso: Fuerza Gravitacional}

\begin{seccion}
\begin{definicion}
    \textbf{Peso ($W$):} Fuerza con la que un planeta atrae a un objeto hacia su centro.
\end{definicion}

\begin{ecuacion}
\textbf{Fórmula del peso:}
$$\vec{W} = m\vec{g}$$

En magnitud:
$$W = mg$$

donde:
\begin{itemize}
    \item $W$ = peso (N)
    \item $m$ = masa (kg)
    \item $g$ = aceleración gravitacional (m/s$^2$)
\end{itemize}
\end{ecuacion}

\textbf{Valores de $g$ en diferentes ubicaciones:}
\begin{itemize}
    \item Tierra (superficie): $g = 9.8$ m/s$^2$ $\approx 10$ m/s$^2$
    \item Luna: $g = 1.6$ m/s$^2$ $\approx g_{Tierra}/6$
    \item Marte: $g = 3.7$ m/s$^2$ $\approx g_{Tierra}/3$
    \item Júpiter: $g = 24.8$ m/s$^2$ $\approx 2.5 g_{Tierra}$
\end{itemize}

\begin{definicion}
    \textbf{Dirección del peso:} Siempre apunta hacia el centro del planeta (verticalmente hacia abajo en la superficie).
\end{definicion}

\end{seccion}

\subsection{Tercera Ley de Newton: Acción y Reacción}

\begin{seccion}
\begin{definicion}
    \textbf{Definición:} Cuando un cuerpo A ejerce una fuerza sobre un cuerpo B, el cuerpo B ejerce simultáneamente una fuerza sobre A de igual magnitud, misma dirección pero sentido opuesto.
\end{definicion}

\begin{ecuacion}
\textbf{Expresión matemática:}
$$\vec{F}_{AB} = -\vec{F}_{BA}$$

donde:
\begin{itemize}
    \item $\vec{F}_{AB}$ = fuerza que A ejerce sobre B (acción)
    \item $\vec{F}_{BA}$ = fuerza que B ejerce sobre A (reacción)
\end{itemize}
\end{ecuacion}

\textbf{Características de los pares acción-reacción:}
\begin{enumerate}
    \item Tienen la misma magnitud
    \item Tienen direcciones opuestas
    \item Actúan sobre cuerpos diferentes
    \item Son simultáneas (ocurren al mismo tiempo)
    \item Son de la misma naturaleza (ambas gravitacionales, ambas de contacto, etc.)
    \item NO se anulan entre sí (actúan sobre objetos distintos)
\end{enumerate}
\newpage
\textbf{Ejemplos de pares acción-reacción:}
\begin{itemize}
    \item Libro sobre mesa: libro empuja mesa hacia abajo $\leftrightarrow$ mesa empuja libro hacia arriba
    \item Caminar: pie empuja suelo hacia atrás $\leftrightarrow$ suelo empuja pie hacia adelante
    \item Cohete: gases expulsados hacia abajo $\leftrightarrow$ cohete empujado hacia arriba
    \item Gravitación: Tierra atrae Luna $\leftrightarrow$ Luna atrae Tierra
\end{itemize}
\end{seccion}

\subsubsection{Diagrama: Tercera Ley de Newton}

\begin{center}
\shorthandoff{>}\shorthandoff{<}
\begin{tikzpicture}[scale=1.2]
    % Libro
    \draw[fill=azulmedio!20, thick] (0,1.5) rectangle (2,2.5);
    \node at (1,2) {\textbf{Libro}};
    
    % Mesa
    \draw[fill=fondogris, thick] (-1,0) rectangle (4,1.5);
    \node at (1.5,0.75) {\textbf{Mesa}};
    
    % Fuerza del libro sobre la mesa (acción)
    \draw[very thick, ->, color=azuloscuro] (1,1.5) -- (1,1.0) node[right, midway] {$\vec{F}_{LM}$};
    
    % Fuerza de la mesa sobre el libro (reacción)
    \draw[very thick, ->, color=verdecorrecto] (1,1.5) -- (1,2.0) node[right, midway] {$\vec{F}_{ML}$};
    
    % Etiquetas
    \node[color=azuloscuro] at (3.5,1.2) {Acción};
    \node[color=verdecorrecto] at (3.5,1.8) {Reacción};
    
    % Relación
    \node at (2, -0.5) {$\vec{F}_{LM} = -\vec{F}_{ML}$};
\end{tikzpicture}
\shorthandon{>}\shorthandon{<}
\end{center}

\section*{Equilibrio Mecánico}

\begin{seccion}
\begin{definicion}
    \textbf{Equilibrio mecánico:} Condición en la que un cuerpo no experimenta cambios en su estado de movimiento (traslacional) ni de rotación.
\end{definicion}

Existen dos tipos de equilibrio:
\begin{enumerate}
    \item \textbf{Equilibrio traslacional:} No hay movimiento de traslación (aceleración lineal = 0)
    \item \textbf{Equilibrio rotacional:} No hay rotación (aceleración angular = 0)
\end{enumerate}
\end{seccion}

\subsection{Primera Condición de Equilibrio: Equilibrio Traslacional}

\begin{seccion}
\textbf{Condición:} Un cuerpo está en equilibrio traslacional cuando la suma vectorial de todas las fuerzas que actúan sobre él es cero.

\begin{ecuacion}
\textbf{Primera condición de equilibrio:}
$$\sum \vec{F} = 0$$

En componentes:
$$\sum F_x = 0 \quad \text{y} \quad \sum F_y = 0$$
\end{ecuacion}


\textbf{Método de resolución de problemas:}
\begin{enumerate}
    \item Dibujar el diagrama de cuerpo libre
    \item Establecer un sistema de coordenadas
    \item Descomponer fuerzas en componentes $x$ y $y$
    \item Aplicar $\sum F_x = 0$ y $\sum F_y = 0$
    \item Resolver el sistema de ecuaciones
\end{enumerate}
\end{seccion}

\subsubsection{Concepto de Torca o Momento de Fuerza}

\begin{seccion}
\begin{definicion}
    \textbf{Torca ($\tau$):} Medida de la efectividad de una fuerza para producir rotación alrededor de un eje o punto.
\end{definicion}

\begin{ecuacion}
\textbf{Definición de torca:}
$$\tau = r F \sin\theta$$

o cuando la fuerza es perpendicular al brazo:
$$\tau = r_\perp F = r F$$

donde:
\begin{itemize}
    \item $\tau$ = torca o momento (N$\cdot$m)
    \item $r$ = distancia del eje de rotación al punto de aplicación (m)
    \item $F$ = magnitud de la fuerza (N)
    \item $\theta$ = ángulo entre $\vec{r}$ y $\vec{F}$
    \item $r_\perp$ = brazo de palanca (distancia perpendicular)
\end{itemize}
\end{ecuacion}

\textbf{Convención de signos (Regla de la mano derecha):}
\begin{itemize}
    \item Torca positiva (+): produce rotación antihoraria 
    \item Torca negativa (-): produce rotación horaria 
\end{itemize}

\textbf{Factores que afectan la torca:}
\begin{enumerate}
    \item \textbf{Magnitud de la fuerza:} Mayor fuerza, mayor torca
    \item \textbf{Brazo de palanca:} Mayor distancia, mayor torca
    \item \textbf{Ángulo de aplicación:} Máxima torca cuando $F \perp r$ ($\theta = 90^\circ$)
\end{enumerate}
\end{seccion}

\subsubsection{Diagrama: Concepto de Torca}

\begin{center}
\shorthandoff{>}\shorthandoff{<}
\begin{tikzpicture}[scale=1.2]
    % Eje de rotación
    \filldraw[color=azuloscuro] (0,0) circle (0.1);
    \draw[thick] (0,-0.3) -- (0,0.3);
    \node[below] at (0,-0.3) {Eje};
    
    % Brazo
    \draw[very thick] (0,0) -- (4,0);
    
    % Punto de aplicación
    \filldraw[color=azulmedio] (4,0) circle (0.08);
    
    % Fuerza perpendicular
    \draw[very thick, ->, color=verdecorrecto] (4,0) -- (4,2.5) node[right, midway] {$\vec{F}$};
    
    % Distancia r
    \draw[<->, thick, color=azuloscuro] (0,-0.5) -- (4,-0.5) node[midway, below] {$r$ (brazo de palanca)};
    
    % Flecha de rotación
    \draw[->, very thick, color=azulmedio] (1,0.3) arc (0:90:1) node[midway, above right] {$\tau = rF$};
    
    \node at (2, -1.5) {Torca máxima: $F$ perpendicular al brazo};
\end{tikzpicture}
\shorthandon{>}\shorthandon{<}
\end{center}

\newpage
\subsection{Segunda Condición de Equilibrio: Equilibrio Rotacional}

\begin{seccion}
\begin{definicion}
    \textbf{Condición:} Un cuerpo está en equilibrio rotacional cuando la suma de todas las torcas que actúan sobre él, respecto a cualquier eje, es cero.
\end{definicion}

\begin{ecuacion}
\textbf{Segunda condición de equilibrio:}
$$\sum \tau = 0$$

o detalladamente:
$$\tau_1 + \tau_2 + \tau_3 + ... = 0$$
\end{ecuacion}

\textbf{Consecuencia:} Si $\sum \tau = 0$, no hay aceleración angular y el cuerpo no rota o rota con velocidad angular constante.

\begin{definicion}
    \textbf{Equilibrio completo:} Un cuerpo está en equilibrio mecánico completo cuando:
$$\sum \vec{F} = 0 \quad \text{y} \quad \sum \tau = 0$$
\end{definicion}

\end{seccion}

\subsection{Aplicación: Balancín y Palancas}

\begin{seccion}
\textbf{Problema típico:} Una barra rígida horizontal está apoyada en un pivote. Diferentes fuerzas actúan a diferentes distancias. ¿Cuál es la condición de equilibrio?

\textbf{Solución:}
\begin{enumerate}
    \item Elegir el eje de rotación (generalmente el pivote)
    \item Calcular la torca de cada fuerza: $\tau_i = r_i F_i$
    \item Asignar signos según la rotación que producirían
    \item Aplicar $\sum \tau = 0$
\end{enumerate}

\textbf{Ventaja mecánica de palancas:} Una fuerza pequeña aplicada a gran distancia puede equilibrar una fuerza grande aplicada a corta distancia:
$$F_1 r_1 = F_2 r_2$$
\end{seccion}

\subsection{Ejemplo de Equilibrio Rotacional}

\begin{center}
\shorthandoff{>}\shorthandoff{<}
\begin{tikzpicture}[scale=1]
    % Barra
    \draw[very thick, fill=fondogris] (-3,0) rectangle (3,0.3);
    
    % Pivote (triángulo)
    \draw[fill=azulmedio, thick] (0,-0.5) -- (-0.3,0) -- (0.3,0) -- cycle;
    
    % Pesos
    \draw[very thick, ->, color=azuloscuro] (-2,0) -- (-2,-1.5) node[right] {$F_1$};
    \draw[very thick, ->, color=azuloscuro] (2,0) -- (2,-1) node[right] {$F_2$};
    
    % Distancias
    \draw[<->, thick] (0,-1.8) -- (-2,-1.8) node[midway, below] {$r_1$};
    \draw[<->, thick] (0,-1.8) -- (2,-1.8) node[midway, below] {$r_2$};
    
    % Condición
    \node at (0, -2.8) {Equilibrio: $F_1 r_1 = F_2 r_2$};
\end{tikzpicture}
\shorthandon{>}\shorthandon{<}
\end{center}

\newpage
\section{Fuerza Elástica: Ley de Hooke}

\begin{seccion}

\begin{defincion}
    \textbf{Ley de Hooke:} Describe la fuerza ejercida por un resorte ideal cuando se estira o comprime desde su posición de equilibrio.
\end{defincion}

\begin{ecuacion}
\textbf{Ecuación de la Ley de Hooke:}
$$F = -kx$$

En magnitud (sin considerar dirección):
$$F = kx$$

donde:
\begin{itemize}
    \item $F$ = fuerza elástica restauradora (N)
    \item $k$ = constante elástica del resorte (N/m)
    \item $x$ = elongación o compresión respecto a la posición natural (m)
    \item El signo negativo indica que la fuerza se opone a la deformación
\end{itemize}
\end{ecuacion}

\textbf{Interpretación física:}
\begin{itemize}
    \item La fuerza es proporcional a la deformación
    \item La fuerza siempre apunta hacia la posición de equilibrio
    \item Mayor $k$ significa resorte más rígido (más difícil de deformar)
    \item Menor $k$ significa resorte más flexible (fácil de deformar)
\end{itemize}
\end{seccion}

\subsection{Constante Elástica}

\begin{seccion}

\begin{defincion}
    \textbf{Constante elástica ($k$):} Característica propia de cada resorte que mide su rigidez.
\end{defincion}

\textbf{Unidad:} N/m (Newton por metro)

\textbf{Significado:}
\begin{itemize}
    \item $k$ grande $\Rightarrow$ resorte rígido (difícil de estirar)
    \item $k$ pequeña $\Rightarrow$ resorte flexible (fácil de estirar)
\end{itemize}

\textbf{Determinación experimental:} Se mide aplicando diferentes fuerzas y midiendo las elongaciones correspondientes:
$$k = \frac{F}{x}$$
\end{seccion}

\newpage
\section{Ley de Gravitación Universal}

\begin{seccion}
\begin{defincion}
    \textbf{Ley de Gravitación Universal de Newton:} Toda partícula en el universo atrae a cualquier otra partícula con una fuerza que es directamente proporcional al producto de sus masas e inversamente proporcional al cuadrado de la distancia que las separa.
\end{defincion}

\begin{ecuacion}
\textbf{Ecuación de la fuerza gravitacional:}
$$F_g = G\frac{m_1 m_2}{r^2}$$

donde:
\begin{itemize}
    \item $F_g$ = fuerza gravitacional (N)
    \item $G$ = constante de gravitación universal = $6.67 \times 10^{-11}$ N$\cdot$m$^2$/kg$^2$
    \item $m_1, m_2$ = masas de los objetos (kg)
    \item $r$ = distancia entre los centros de masa (m)
\end{itemize}
\end{ecuacion}

\textbf{Características de la fuerza gravitacional:}
\begin{enumerate}
    \item Actúa a distancia (no requiere contacto)
    \item Es una fuerza mutua (par acción-reacción)
    \item Es muy débil comparada con otras fuerzas fundamentales
\end{enumerate}
\end{seccion}

\subsection{Relación entre Peso y Gravitación Universal}

\begin{seccion}
El peso de un objeto en la superficie de un planeta es un caso particular de la ley de gravitación universal.

Para un objeto de masa $m$ en la superficie de la Tierra (masa $M_T$, radio $R_T$):

\begin{ecuacion}
$$W = G\frac{M_T m}{R_T^2} = m \left(G\frac{M_T}{R_T^2}\right) = mg$$

donde:
$$g = G\frac{M_T}{R_T^2} \approx 9.8 \text{ m/s}^2$$
\end{ecuacion}


\end{seccion}

\subsection{Variación de la Fuerza Gravitacional con la Distancia}

\begin{seccion}
\textbf{Dependencia cuadrática inversa:} La fuerza disminuye rápidamente al aumentar la distancia.

\textbf{Ejemplos:}
\begin{itemize}
    \item Si la distancia se duplica ($r' = 2r$): $F' = F/4$ (se reduce a la cuarta parte)
    \item Si la distancia se triplica ($r' = 3r$): $F' = F/9$ (se reduce a la novena parte)
    \item Si la distancia se reduce a la mitad ($r' = r/2$): $F' = 4F$ (se cuadruplica)
\end{itemize}

\end{seccion}

\subsubsection*{Gráfica: Fuerza Gravitacional vs Distancia}

\begin{center}
\shorthandoff{>}\shorthandoff{<}
\begin{tikzpicture}
\begin{axis}[
    width=11cm,
    height=7cm,
    axis lines=middle,
    xlabel={Distancia $r$ (unidades arbitrarias)},
    ylabel={Fuerza $F_g$ (unidades arbitrarias)},
    xmin=0, xmax=5,
    ymin=0, ymax=10,
    grid=major,
    thick,
    every axis x label/.style={at={(current axis.right of origin)},anchor=west},
    every axis y label/.style={at={(current axis.above origin)},anchor=south},
    domain=0.5:5,
    samples=100
]

% Curva 1/r²
\addplot[color=azuloscuro, very thick] {4/x^2};

% Puntos de referencia
\addplot[only marks, mark=*, color=verdecorrecto, mark size=3pt] coordinates {(1, 4) (2, 1) (4, 0.25)};

% Etiquetas
\node[color=azuloscuro] at (axis cs: 3.5, 6) {$F_g \propto \frac{1}{r^2}$};

\end{axis}
\end{tikzpicture}
\shorthandon{>}\shorthandon{<}
\end{center}

\section{Movimiento Planetario: Leyes de Kepler}

\begin{seccion}
Johannes Kepler (1571-1630) formuló tres leyes empíricas que describen el movimiento de los planetas alrededor del Sol. Posteriormente, Newton demostró que estas leyes se derivan de su Ley de Gravitación Universal.
\end{seccion}

\subsection{Primera Ley de Kepler: Ley de las Órbitas}

\begin{seccion}
\textbf{1ra Ley de Kepler:} Los planetas describen órbitas elípticas alrededor del Sol, con el Sol en uno de los focos de la elipse.

\begin{ecuacion}
\textbf{Elementos de la órbita elíptica:}
\begin{itemize}
    \item \textbf{Afelio:} Punto más alejado del Sol
    \item \textbf{Perihelio:} Punto más cercano al Sol
    \item \textbf{Semieje mayor ($a$):} Mitad de la distancia entre afelio y perihelio
\end{itemize}
\end{ecuacion}

\begin{nota}
    \textbf{Observación importante:} Aunque las órbitas son elípticas, la mayoría de los planetas tienen órbitas casi circulares (baja excentricidad).
\end{nota}

\end{seccion}

\subsubsection{Diagrama: Órbita Elíptica}

\begin{center}
\shorthandoff{>}\shorthandoff{<}
\begin{tikzpicture}[scale=1.2]
    % Elipse (órbita)
    \draw[very thick, color=azulmedio] (0,0) ellipse (4 and 2.5);
    
    % Focos
    \filldraw[color=azuloscuro] (-2.65,0) circle (0.12);
    \node[below, color=azuloscuro] at (-2.65,-0.2) {Sol (foco)};
    \filldraw[color=fondogris] (2.65,0) circle (0.08);
    
    % Planeta
    \filldraw[color=verdecorrecto] (4,0) circle (0.1);
    \node[right] at (4,0) {Planeta};
    
    % Perihelio y Afelio
    \draw[<->, thick] (-4,0) -- (4,0);
    \node[below] at (-4,-0.2) {Perihelio};
    \node[below] at (4,-0.2) {Afelio};
    
    % Semieje mayor
    \draw[<->, color=azuloscuro, thick] (0,-3) -- (4,-3);
    \node[below, color=azuloscuro] at (2,-3) {$a$ (semieje mayor)};
\end{tikzpicture}
\shorthandon{>}\shorthandon{<}
\end{center}

\subsection{Segunda Ley de Kepler: Ley de las Áreas}

\begin{seccion}
\begin{defincion}
    \textbf{2da Ley de Kepler::} El vector posición que une el Sol con un planeta barre áreas iguales en tiempos iguales.
\end{defincion}

\textbf{Consecuencia:} Los planetas se mueven más rápido cuando están cerca del Sol (perihelio) y más lento cuando están lejos (afelio).

\textbf{Principio físico:} Esta ley refleja la conservación del momento angular del sistema.

\begin{ecuacion}
\textbf{Velocidad areolar constante:}
$$\frac{dA}{dt} = \text{constante}$$

donde $A$ es el área barrida y $t$ es el tiempo.
\end{ecuacion}
\end{seccion}
\newpage
\subsection{Tercera Ley de Kepler: Ley de los Períodos}

\begin{seccion}

\textbf{3ra Ley de Kepler:} El cuadrado del período orbital de un planeta es directamente proporcional al cubo del semieje mayor de su órbita.

\begin{ecuacion}
\textbf{Tercera Ley de Kepler:}
$$T^2 = k a^3$$

o en forma de proporción:
$$\frac{T^2}{a^3} = \text{constante (para todos los planetas)}$$

Para dos planetas diferentes:
$$\frac{T_1^2}{T_2^2} = \frac{a_1^3}{a_2^3}$$

donde:
\begin{itemize}
    \item $T$ = período orbital (tiempo para dar una vuelta completa)
    \item $a$ = semieje mayor de la órbita
    \item $k$ = constante que depende de la masa del Sol
\end{itemize}
\end{ecuacion}

\textbf{Aplicación práctica:} Permite calcular el período de un planeta conociendo su distancia al Sol, o viceversa.
\end{seccion}


\newpage
\section{Formulario}

\begin{nota}
\textbf{Componentes Vectoriales:}
$$F_x = F\cos\theta \quad , \quad F_y = F\sin\theta$$
$$F = \sqrt{F_x^2 + F_y^2}$$

\vspace{0.3cm}

\textbf{Leyes de Newton:}
\begin{itemize}
    \item \textbf{Primera Ley:} $\sum \vec{F} = 0 \Rightarrow \vec{v} = \text{constante}$
    \item \textbf{Segunda Ley:} $\vec{F}_{neta} = m\vec{a}$
    \item \textbf{Tercera Ley:} $\vec{F}_{AB} = -\vec{F}_{BA}$
\end{itemize}

\vspace{0.3cm}

\textbf{Peso:}
$$W = mg$$

\vspace{0.3cm}

\textbf{Equilibrio:}
\begin{itemize}
    \item \textbf{Traslacional:} $\sum \vec{F} = 0$
    \item \textbf{Rotacional:} $\sum \tau = 0$, donde $\tau = rF\sin\theta$
\end{itemize}

\vspace{0.3cm}

\textbf{Ley de Hooke:}
$$F = kx$$

\vspace{0.3cm}

\textbf{Gravitación Universal:}
$$F_g = G\frac{m_1 m_2}{r^2}$$

\vspace{0.3cm}

\textbf{Tercera Ley de Kepler:}
$$\frac{T_1^2}{T_2^2} = \frac{a_1^3}{a_2^3}$$
\end{nota}

